% GENERATED FILE. Edit canonical lessons, then run npm run generate:books.
% canonical-source-hash: fnv1a64:373a526aa186224d
% canonical-lessons: TA-C45-eye, TA-C45-ear, TA-C45-nose, TA-C45-tooth, TA-C45-stomach

\chapter{More of the Body}
\label{ch:ta-body-more}
\hlchaptermodality{45}

\begin{chapteropening}
\textbf{By the end of this chapter:} \emph{I can name five more parts of the body in Tamil, and point to the sound changes that separate them from their Malayalam, Telugu and Kannada cousins.}

Complete TA-C45-stomach and say all five words in order.
\end{chapteropening}

\section[kaṇ]{\ta{கண்} (kaṇ) — an eye}

\label{lesson:TA-C45-eye}



\begin{quote}
Before the new run of words: what did \emph{uppu} mean, and what did \emph{puttakam} mean?
\end{quote}

\subsection*{You'll want to know: \ta{கண்}}
\textbf{\ta{கண்}} (\emph{kaṇ}) — \textquotedblleft{}an eye\textquotedblright{}.

Proto-Dravidian \textbf{*kaṇ}, held almost still everywhere it went: Malayalam \ml{കണ്ണ്}, Telugu \te{కన్ను}, Kannada \ka{ಕಣ್ಣು}. Tamil's is the shortest of the four, one syllable and done.

It reaches well past the face. A \emph{kaṇ} is the joint on a stalk of sugarcane and the mesh of a net — wherever a thing has an opening or a knot, Tamil is willing to call it an eye.

\emph{talai} and \emph{kai}, the head and the hand, came earlier. The face gets its own turn now, one piece at a time.

\begin{grammarlens}[title={What you've built}]
The first of five more parts of the body.
\end{grammarlens}

\subsection*{Your turn}
Say these aloud:

\begin{itemize}
\raggedright
  \item \emph{kaṇ}
  \item it once more, slowly
  \item \emph{kaṇ}, then \emph{talai}, so the part and the whole sit together
\end{itemize}

\begin{tcolorbox}[breakable,colback=teal!4,colframe=teal!35!black,title={Before you move on}]
What does \ta{கண்} mean? (\textquotedblleft{}an eye\textquotedblright{}.) Say it once more.
\end{tcolorbox}
\section[kādu]{\ta{காது} (kādu) — an ear}

\label{lesson:TA-C45-ear}



\begin{quote}
Before the new one: what did \emph{kaṇ} mean?
\end{quote}

\subsection*{You'll want to know: \ta{காது}}
\textbf{\ta{காது}} (\emph{kādu}) — \textquotedblleft{}an ear\textquotedblright{}.

Here it is Tamil that stands apart. \emph{kādu} is the ear a Tamil speaker has, and the other three sisters built their everyday word on a different root altogether: Malayalam \ml{ചെവി}, Telugu \te{చెవి}, Kannada \ka{ಕಿವಿ}.

Tamil owns that root as well. \ta{செவி} (\emph{sevi}) names the same organ and is kept for poetry and formal speech — so the word the sisters use over breakfast is the one Tamil holds back for its high register. Two words, one ear, and the occasion decides which.

\begin{grammarlens}[title={What you've built}]
Two: an eye and an ear.
\end{grammarlens}

\subsection*{Your turn}
Say these aloud:

\begin{itemize}
\raggedright
  \item \emph{kādu}
  \item it once more, slowly
  \item \emph{kaṇ}, then \emph{kādu}, and hear the long \emph{ā} arrive in the second
\end{itemize}

\begin{tcolorbox}[breakable,colback=teal!4,colframe=teal!35!black,title={Before you move on}]
What does \ta{காது} mean? (\textquotedblleft{}an ear\textquotedblright{}.) Say it once more.
\end{tcolorbox}
\section[mūkku]{\ta{மூக்கு} (mūkku) — a nose}

\label{lesson:TA-C45-nose}



\begin{quote}
Before the new one: what did \emph{kādu} mean?
\end{quote}

\subsection*{You'll want to know: \ta{மூக்கு}}
\textbf{\ta{மூக்கு}} (\emph{mūkku}) — \textquotedblleft{}a nose\textquotedblright{}.

Another word the family barely touched: Malayalam \ml{മൂക്ക്}, Telugu \te{ముక్కు}, Kannada \ka{ಮೂಗು}. Only the ending moves.

\emph{mūkku} also names the point of a thing — the prow of a boat, the tip of a pen. Whatever sticks out in front is the nose of it.

\begin{grammarlens}[title={What you've built}]
Three, and all three are on the face.
\end{grammarlens}

\subsection*{Your turn}
Say these aloud:

\begin{itemize}
\raggedright
  \item \emph{mūkku}
  \item it once more, slowly
  \item \emph{kādu}, then \emph{mūkku}, and let the second start long
\end{itemize}

\begin{tcolorbox}[breakable,colback=teal!4,colframe=teal!35!black,title={Before you move on}]
What does \ta{மூக்கு} mean? (\textquotedblleft{}a nose\textquotedblright{}.) Say it once more.
\end{tcolorbox}
\section[pal]{\ta{பல்} (pal) — a tooth}

\label{lesson:TA-C45-tooth}



\begin{quote}
Before the new one: what did \emph{mūkku} mean?
\end{quote}

\subsection*{You'll want to know: \ta{பல்}}
\textbf{\ta{பல்}} (\emph{pal}) — \textquotedblleft{}a tooth\textquotedblright{}.

Proto-Dravidian \textbf{*pal}, and the sisters pulled it three ways. Malayalam \ml{പല്ല്}, Kannada \ka{ಹಲ್ಲು}, Telugu \te{పన్ను}. Kannada is the odd one again, with the same \emph{p-} to \emph{h-} shift the fruit-word showed; Telugu turned the \emph{l} into an \emph{n}. Tamil stops the word dead on its \emph{l} and adds nothing.

In everyday speech it often comes out \emph{pallu}, with the \emph{l} doubled — which is exactly the shape Malayalam and Kannada put on the page.

\begin{grammarlens}[title={What you've built}]
Four, and the first one you can bite with.
\end{grammarlens}

\subsection*{Your turn}
Say these aloud:

\begin{itemize}
\raggedright
  \item \emph{pal}
  \item it once more, slowly
  \item \emph{mūkku}, then \emph{pal}, and cut the second one short
\end{itemize}

\begin{tcolorbox}[breakable,colback=teal!4,colframe=teal!35!black,title={Before you move on}]
What does \ta{பல்} mean? (\textquotedblleft{}a tooth\textquotedblright{}.) Say it once more.
\end{tcolorbox}
\section[vayiṟu]{\ta{வயிறு} (vayiṟu) — the stomach}

\label{lesson:TA-C45-stomach}



\begin{quote}
Before the new one: what did \emph{pal} mean?
\end{quote}

\subsection*{You'll want to know: \ta{வயிறு}}
\textbf{\ta{வயிறு}} (\emph{vayiṟu}) — \textquotedblleft{}the stomach\textquotedblright{}.

The \ta{ற} in the middle is the hard \emph{ṟ} this book met at \emph{aṟai}, 'a room' — not the soft \emph{r} of \emph{peyar}. Malayalam says \ml{വയറ്}, the same word; Telugu and Kannada went elsewhere entirely for this one.

Where an English speaker says they are hungry, Tamil is happy to talk about the \emph{vayiṟu} and leave the person out of it: the belly is what is empty, and you merely own it. That is the same habit behind \emph{uṅgaḷ uḍambu eppaḍi irukku}, where the health question hangs on the body rather than on you.

\begin{grammarlens}[title={What you've built}]
That closes the run: an eye, an ear, a nose, a tooth, and what holds the food.
\end{grammarlens}

\subsection*{Your turn}
Say these aloud:

\begin{itemize}
\raggedright
  \item \emph{vayiṟu}
  \item it once more, slowly
  \item all five in order — \emph{kaṇ}, \emph{kādu}, \emph{mūkku}, \emph{pal}, \emph{vayiṟu}
\end{itemize}

\begin{tcolorbox}[breakable,colback=teal!4,colframe=teal!35!black,title={Before you move on}]
What does \ta{வயிறு} mean? (\textquotedblleft{}the stomach\textquotedblright{}.) Say it once more.
\end{tcolorbox}
